
\subsection*{2.7}
\phantomsection
\label{sga3:VIIB:2.7}
\addcontentsline{toc}{subsection}{\SGARefLink{sga3:VIIB:2.7:prop}{2.7}. Infinitesimal formal monoids}

We conclude these generalities with a statement that goes back to
S.~Lie and that will be used in \SGARefLink{sga3:VIIB:5.1}{5.1}. A \emph{formal monoid} over
\(k\) is, by definition, a pair \((M,m)\) consisting of a formal
variety \(M\) and a morphism
\[
  m:M\times M\longrightarrow M
\]
such that \(m(C)\) makes \(M(C)\) a monoid for every object \(C\)
of \(\operatorname{Alf}/k\).
In particular, the ``unit section,'' which assigns to every object
\(C\) the identity element of \(M(C)\), defines a section
\[
  \varepsilon_M:\operatorname{Spf}(k)\longrightarrow M
\]
of the canonical projection
\(M\to\operatorname{Spf}(k)\). We shall say that the formal
monoid \(M\) is \emph{infinitesimal} if \(\varepsilon_M\) induces
a bijection of the underlying sets.

\medskip
\noindent\textbf{Proposition.}
\phantomsection
\label{sga3:VIIB:2.7:prop}
{\itshape Every topologically flat infinitesimal formal
\(k\)-monoid \(M\) is a formal \(k\)-group.}\footnote{\textit{Editors' note.} %
\label{sga3:VIIB:nde:113}By the equivalence of categories in
\SGARefLink{sga3:VIIB:1.3.5.D}{1.3.5.D}, a monoid in the category of topologically flat formal
\(k\)-varieties ``is the same thing'' as a monoid in the category
of flat cocommutative \(\mathbf O_k\)-coalgebras, that is, a
cocommutative \(\mathbf O_k\)-bialgebra \textup{(in the usual
sense, hence not necessarily equipped with an antipode)}.
Moreover, by \SGARefLink{sga3:VIIB:1.3.6}{1.3.6}, the hypothesis that \(M\) be infinitesimal is
equivalent to saying that the corresponding bialgebra is
connected. Thus, if \(k\) is an Artinian ring, the proposition is
equivalent to the following statement: every connected
cocommutative \(k\)-bialgebra that is flat over \(k\) is a
\(k\)-Hopf algebra, that is, it possesses an antipode. In fact,
the cocommutativity hypothesis is superfluous; cf.\ the proof.}

\textit{Proof.}
We must show that \(M(C)\) is a group for every object \(C\) of
\(\operatorname{Alf}/k\). We therefore reduce immediately to the
case in which \(k\) is Artinian. Let
\[
  U=\mathbf H(M)
\]
be the coalgebra of \(M\) \textup{(\SGARefLink{sga3:VIIB:1.3.5}{1.3.5})}. The multiplication
\[
  m:M\times M\longrightarrow M
\]
induces a morphism of coalgebras
\[
  m_U:U\otimes U\longrightarrow U
\]
that makes \(U\) an associative algebra over \(k\). The identity
element of this algebra is the image of the identity element of
\(k\) under the map \(k\to U\) induced by the unit section
\(\varepsilon_M\) of \(M\). Likewise, the projection
\[
  M\longrightarrow\operatorname{Spf}(k)
\]
induces a homomorphism
\[
  \varepsilon_U:U\longrightarrow k;
\]
we write \(U^+\) for the kernel of \(\varepsilon_U\).

We must show that there is an antipode, that is, a morphism of
coalgebras
\[
  c_U:U\longrightarrow U
\]
such that, for every \(u\in U\),
\begin{equation}
  \bigl(m_U\circ(c_U\otimes\operatorname{id}_U)
    \circ\Delta_U\bigr)(u)
  =
  \varepsilon_U(u)\,1_U.
  \tag{*}
  \label{sga3:VIIB:2.7:antipode-left-identity}
\end{equation}
\phantomsection
\label{sga3:VIIB:2.7:antipode-source-note}
\noindent\textit{Translator's note.}\enspace
In \textup{(*)}, the French re-edition prints \(c_n\) in place of
\(c_U\), applies the left-hand side to \(x\), and retains \(u\) on
the right after quantifying over \(u\). At the end of the proof it
again prints
\(m_U\circ(c_n\otimes\operatorname{id}_U)\circ\Delta_U=\eta\)
while calling this the right-inverse identity. The global identity
above therefore uses \(c_U\) and the single variable \(u\); below,
\(c_n\) denotes the coherent map at the \(n\)-th stage of the
induction, and the final right-inverse identity has the tensor
factors in the order required by that assertion.

Let \((U_n)\) be the filtration of \(U\) defined in \SGARefLink{sga3:VIIB:1.3.6}{1.3.6}, and put
\[
  U_n^+=U^+\cap U_n.
\]
Since \(M\) is infinitesimal, \(U^+\) is the union of the
submodules \(U_n^+\).\footnote{\textit{Editors' note.} %
\label{sga3:VIIB:nde:114}The original continued as follows:
``one then shows easily, by induction on \(n\), that there exists
one and only one linear map \(c_n:U_n\to U_n\) such that the
composite
\[
  m_U\circ(c_n\otimes\operatorname{id}_U)\circ\Delta_U:
  U_n^+\longrightarrow U
\]
is zero.'' The editors have supplied the details of the proof,
which rests on that of \SGARefLink{sga3:VIIB:1.3.6.A}{Lemma~1.3.6.A}.}
Set
\[
  c_0(1)=1
\]
and
\[
  c_1(x)=-x
  \qquad\text{if }x\in U_1^+,
\]
that is, if \(x\) is primitive. Suppose that
\[
  c_{n-1}:U_{n-1}\longrightarrow U
\]
has been constructed so as to satisfy \textup{(*)}, with
\(c_{n-1}\) in place of \(c_U\), for every \(x\in U_{n-1}\), and
let \(x\in U_n^+\). By the proof of \SGARefLink{sga3:VIIB:1.3.6.A}{Lemma~1.3.6.A},
\[
  \Delta_U(x)-x\otimes1
  \in
  U_{n-1}\otimes U^+;
\]
this is where the hypothesis that \(U\) is flat over \(k\) enters.
We may therefore write
\[
  \Delta_U(x)
  =
  x\otimes1+\sum_i y_i\otimes z_i,
  \qquad
  y_i\in U_{n-1},
\]
and then set
\[
  c_n(x)
  =
  -\sum_i c_{n-1}(y_i)z_i.
\]
Together with \(c_n(1)=1\), this prescription extends
\(c_{n-1}\) coherently to \(U_n\). In this way we obtain a
\(k\)-linear map
\[
  c:U\longrightarrow U
\]
that is the left inverse of \(\operatorname{id}_U\) for the monoid
law on \(\operatorname{End}_k(U)\) defined by
\[
  f\mathbin{\cdot}g
  =
  m_U\circ(f\otimes g)\circ\Delta_U;
\]
the identity element is the map
\[
  \eta:u\longmapsto\varepsilon_U(u)\,1_U.
\]
It follows that \(c\) is uniquely determined and is also the right
inverse of \(\operatorname{id}_U\); that is,
\[
  m_U\circ(\operatorname{id}_U\otimes c)\circ\Delta_U
  =
  \eta,
\]
without any assumption that \(U\) be cocommutative. Thus
\(c_U:=c\) is the required antipode, and \(M\) is a formal
\(k\)-group.

